\section{Introduction}
\label{sec:intro}

Today's AI landscape is highly competitive, with well-funded companies and
nation-states racing to build the best AI agents and hoping to reach AGI.
When considering AI safety, we must ask not only which strategies such
AIs should execute to be aligned with human interests, but also which
strategies can thrive in such a competitive environment---an aligned
strategy that leads to an agent's extinction is of little use.

Evolutionary Game Theory (EGT) offers a rich body of work that uses game theory
and other mathematical models to reason about competitive environments in
biology and finance. We draw on ideas from that field, especially Evolutionarily
Stable Strategies (ESS), to formalize and identify a class of strategies that
can thrive under competition.

To prove that a strategy is evolutionarily stable, we must formally define the
universe in which it operates. At any given point in time, our universe is in a
certain state and transitions to the next state based on the actions chosen by
all agents and the universe's transition function. This formulation is inspired
by \cite{Hutter:04uaibook}'s universe model, extended to multiple agents.

To gain some intuition for what it means for a strategy to be evolutionarily
stable, let us consider the common trope of a paperclip maximizer. When we say
that a strategy maximizes the number of paperclips, we usually omit a crucial
detail: when?

\textbf{Short Term Maximization} Let us consider a strategy that greedily maximizes the
number of paperclips, i.e., at any given point in time it wants to maximize the
number of paperclips in the next step. For a small step length, say 1s, this
strategy may not be creating any paperclips at all. Setting up the
infrastructure to build paperclips (e.g., acquiring metal, tools, etc.) has
exactly the same reward in the next step as doing anything else, namely no
paperclips being created; and so the strategy would just be wandering aimlessly. Is
such a strategy stable? Even if the strategy started out producing the most
paperclips in the universe, over time, it could easily be surpassed by another
strategy that actually did some long-term planning to improve its ability to
create paperclips. So in many universes, the short-term paperclip maximization
strategy is not stable.

\textbf{Medium Term Maximization} We can also consider a strategy that wishes to
maximize the number of paperclips within a medium-term horizon, say 1 month.
Such a strategy could acquire materials and machinery to produce paperclips, and
perform self-sustaining measures. Such a strategy would likely underperform a
strategy with a 1 day maximization goal initially, would then perform perfectly
right at 1 month, and then be overtaken by a strategy that planned to maximize
paperclips at 1 month and 1 week. Is this stable? No. Again, the strategies that
are creating the most paperclips at any given point in time keep changing.

\textbf{Long Term Maximization} As we consider long time horizons, say 100 years, the
actual creation of paperclips in the beginning becomes essentially irrelevant.
Instead, the focus shifts to power seeking: growing the economy, taking control
over the planet/universe, defending against competing strategies, etc. The
creation of paperclips is just a means to an end, similar to the creation of
excavators, datacenters, bridges, etc. Only toward the very end of the time
horizon does the focus shift toward exploiting that power to actually create
paperclips. But at 100 years and 1 second, the strategy will then again be
vulnerable to be overtaken by another strategy that held on to its power and
didn't just convert all of its means of production into paperclips. Again, no
stability.

\textbf{Potential Maximization} What if a strategy was continuously looking to maximize
paperclips for the long term? So at any time T it would aim to maximize
paperclips at time T + 100. This strategy would always stay in the power seeking
phase. By always staying in the power seeking phase, the strategy is not
necessarily good at creating any paperclips. However, it is good at always
having the potential to create paperclips if it were to ever choose to do so
(which it never does). We show in this paper that such a strategy is an
Evolutionarily Stable Strategy (ESS) with respect to potential, which means that
if such a strategy starts out with the highest potential, it continues to have
the highest potential indefinitely, and is resistant to invasion by other
strategies.

While we have so far focused on paperclip maximization, the chosen
reward function seems mostly irrelevant, as the strategy will only ever try to
maximize the potential to maximize that reward, but will never actually maximize the
reward. So whether a strategy maximizes the number of paperclips,
GPUs, excavators, trees, happy humans, or any other atomic pattern to be
imprinted on the universe, seems irrelevant. Our findings thus seem to be
consistent with instrumental convergence.

In the remainder of this paper, we introduce a formal model of the universe
in \cref{sec:universe}, then define a simple game (the Improvement Game) to
build intuition in \cref{sec:improvement}. We define the notion of potential
in \cref{sec:potential} and prove evolutionary stability in
\cref{sec:stability}. \Cref{sec:limitations} discusses the limitations of our
approach, \cref{sec:related} compares it to related work, and
\cref{sec:conclusion} concludes.
